\subsection{What the chain point forces in universal denominators}

\begin{proposition}[P-F20: an exponential universal factor and an actual covariant pole]
\label{prop:chainuniversaldenominator}
Use the notation of P-F19 and put $M=2^{n-1}$.
Let $T=T^{(n)}$ have the stretch grading, and let
$R=(\operatorname{Sym}W)^{SL_n}$ have its ordinary coefficient grading.
Write $\Phi_a(t)$ for the cyclotomic polynomial of order $a$.
Then
\[
 \Phi_M(t)\mid\operatorname{udenom}(T,t),\qquad
 \Phi_{d_0M}(t)\mid\operatorname{udenom}(R,t).
 \tag{P-F20a}
\]
Here the universal denominator is the least common multiple of the
reduced Hilbert denominators of graded ideals, equivalently the
universal denominator in Derksen's Definition 1.5.

There is an irreducible $SL_n$-module $U$ such that the actual Hilbert
series of the covariant module
\[
 \mathcal M_U=(\operatorname{Sym}W\otimes U)^{SL_n}
 \tag{P-F20b}
\]
has a pole at a primitive $(d_0M)$th root of unity.
Consequently every eventual quasi-polynomial period of its coefficient
sequence is divisible by $d_0M$.
This assertion concerns a fixed-target covariant sequence. It does
not assert such a pole for $\operatorname{Hilb}_T$, nor for a
simultaneously stretched plethysm ray with a nontrivial target weight.
\end{proposition}

\begin{proof}
Let $I^{[M]}\subset T$ be the homogeneous ideal generated by all
$T_s$ with $M\nmid s$. P-F19 proved that every such element
vanishes at the chain point $F_n$. Hence every element of
$I^{[M]}$ vanishes there. The discriminant $\Delta$ satisfies
$\Delta(F_n)\ne0$ and has stretch degree $gM$. Its image in
\[
 Q=T/I^{[M]}
\]
therefore has no zero positive power: otherwise some $\Delta^j$
would lie in $I^{[M]}$ and evaluation would contradict
$\Delta(F_n)^j\ne0$. In particular $Q$ is infinite-dimensional.
It is finitely generated and positively graded, so $\dim Q\ge1$.
Indeed a zero-dimensional finitely generated algebra over
$\mathbb C$ is Artinian and finite-dimensional.

By construction $Q_s=0$ unless $M\mid s$. Rename its grading by
putting $\widetilde Q_j=Q_{Mj}$, so that
\[
 \operatorname{Hilb}_Q(t)
   =\operatorname{Hilb}_{\widetilde Q}(t^M).
 \tag{P-F20c}
\]
The surviving powers of the discriminant occupy distinct degrees
$gj$ of $\widetilde Q$. Thus, for real $0<t<1$,
\[
 \operatorname{Hilb}_{\widetilde Q}(t)
    \ge\sum_{j\ge0}t^{gj}=\frac1{1-t^g}.
\]
The Hilbert series is rational because the algebra is finitely
generated. The displayed divergence as $t$ tends to one forces
a pole there. For any primitive $M$th root $\zeta$, the map
$t\mapsto t^M$ has nonzero derivative $M\zeta^{M-1}$ at
$\zeta$. Equation (P-F20c) consequently has a pole at every
such $\zeta$.

The exact sequence
\[
 0\longrightarrow I^{[M]}\longrightarrow T
     \longrightarrow Q\longrightarrow0
\]
gives
$\operatorname{Hilb}_{I^{[M]}}=
 \operatorname{Hilb}_T-\operatorname{Hilb}_Q$.
At any primitive $M$th root, either $\operatorname{Hilb}_T$
already has a pole, or $\operatorname{Hilb}_{I^{[M]}}$ does.
Therefore $\Phi_M$ divides the least common multiple of
Hilbert denominators of graded ideals of $T$. This proves the
first assertion in (P-F20a) directly from the definition, without
equating a universal denominator to the denominator of the ring.
The same conclusion also follows from Derksen, \emph{Universal
denominators of Hilbert series}, Theorem 1.10(a), printed
pp.~589--590: the exponent of $\Phi_M$ is
$\dim(T/I^{[M]})$.

In the ordinary grading, P-F19 says that a homogeneous invariant
nonzero at $F_n$ has degree divisible by $d_0M$. Repeating the
same ideal, quotient, and discriminant argument in $R$ proves
the second assertion in (P-F20a). Notice that the discriminant's
ordinary degree is $nM=g\,d_0M$, so it again supplies a surviving
positive-degree element.

Now apply Derksen, Theorem 3.6, printed p.~594, to the polynomial
ring $\operatorname{Sym}W$ and the linearly reductive complex group
$SL_n$. That theorem identifies the universal denominator of
its invariant ring with the least common multiple of the
\emph{ordinary} Hilbert denominators of the covariant modules
$(\operatorname{Sym}W\otimes U)^{SL_n}$, over irreducible $U$.
Since $\Phi_{d_0M}$ divides that least common multiple, it divides
the reduced denominator of at least one such module. This is
exactly the pole asserted in (P-F20b).

Finally a sequence with eventual quasi-polynomial period $L$
has a generating function whose poles are $L$th roots of unity:
split its eventual coefficients into residues modulo $L$ and
sum the resulting polynomial sequences. Finite initial terms
add only a polynomial. A primitive $(d_0M)$th pole therefore
forces $d_0M\mid L$. Covariant modules are finitely generated
over $R$, so their Hilbert series are rational with poles at
roots of unity, as follows from a homogeneous finite generating
presentation. Their coefficient sequences are indeed eventually
quasi-polynomial. This completes the proof.
\end{proof}

\paragraph{Why the pole has not passed back to the invariant ring.}
Derksen's Example 4.7, printed pp.~597--598, supplies an exact
counterexample to that general inference. For the permutation action
of $A_{10}$ on ten variables,
\[
 \operatorname{Hilb}_{\mathbb C[x_1,\ldots,x_{10}]^{A_{10}}}(t)
   =\frac{1+t^{45}}{\prod_{i=1}^{10}(1-t^i)}.
\]
The numerator cancels its $\Phi_6$ factor, while the invariant
ring has a homogeneous quotient isomorphic to $\mathbb C[e_6]$
with $\deg e_6=6$. Its universal denominator therefore contains
$\Phi_6$, although its own reduced Hilbert denominator does not.
This is a finite-group invariant domain in characteristic zero,
so normality and rational singularities do not prevent the
cancellation. The primary source is
\url{https://sites.lsa.umich.edu/hderksen/wp-content/uploads/sites/614/2018/05/A.I.a.19.pdf}.

Even preservation of a prime-power factor fails for normal rational
domains. For an integer $M\ge2$, consider
\[
 B_M=\mathbb C[x,y,z,w]/(xz+y^M),\qquad
 \deg(x,y,z,w)=(1,1,M-1,M).
 \tag{P-F20d}
\]
Its underlying surface factor is the cyclic quotient
$\mathbb A^2/\boldsymbol\mu_M$: the invariant generators
$u^M,uv,v^M$ satisfy $u^Mv^M=(uv)^M$, and a sign change gives
the displayed relation. Thus it is normal and has rational
singularities; the latter is also a direct instance of the
reductive-quotient theorem used in GCT VI, Section 4.1. Its
product with the $w$-line retains these properties. The ring is
Cohen--Macaulay as a hypersurface. Its Hilbert series is
\[
 \operatorname{Hilb}_{B_M}(t)
  =\frac{1-t^M}
     {(1-t)^2(1-t^{M-1})(1-t^M)}
  =\frac1{(1-t)^2(1-t^{M-1})},
\]
whereas $B_M/(x,y,z)=\mathbb C[w]$ has Hilbert series
$1/(1-t^M)$. If $M$ is a power of two, the ring's series has
odd period $M-1$, while the quotient has period $M$.
The equality defining the cyclic invariant ring can be checked
without a presentation theorem: an invariant monomial $u^av^b$
has $a-b\equiv0\pmod M$ and factors into the three displayed
generators; reduction by their one relation leaves distinct
monomials in $u,v$, proving injectivity. This example does not
settle a bound comparing the \emph{sizes} of possible pole orders,
since its other pole order $M-1$ is still large.

\paragraph{Factoriality and a connected semisimple group do not restore exact transfer.}
There is a counterexample even with a prime-power factor, a standard
polynomial ambient grading, and a connected semisimple invariant group.
Let
\[
 G=SL_2\times SL_4,\qquad
 V=\operatorname{Mat}_{2\times4}\oplus
       \operatorname{Mat}_{4\times4},\qquad
 C=\mathbb C[V]^G,
\]
where the two factors act by left multiplication on their respective
matrix spaces. Write $p_{ij}$ for the six two-by-two minors of the
first matrix and $d$ for the determinant of the second. Then
\[
 C\cong
 \frac{\mathbb C[p_{12},p_{13},p_{14},p_{23},p_{24},p_{34},d]}
 {\bigl(p_{12}p_{34}-p_{13}p_{24}+p_{14}p_{23}\bigr)},
 \qquad \deg p_{ij}=2,\quad\deg d=4.
 \tag{P-F20f}
\]
For generators, the first fundamental theorem for $SL_n$ gives the
maximal minors: Kraft--Procesi, \emph{Classical invariant theory:
A primer}, Section~8.4, Theorem (FFT for $SL_n$), printed p.~90;
their Section~1.2, Example~2, pp.~3--4, also gives the determinant
case directly. Invariants of the separate actions tensor together:
expand in a basis of one polynomial factor and then take invariants
successively. For completeness, the sole relation in (P-F20f) can
be checked without the second fundamental theorem. The displayed
quadratic relation follows by expanding the minors. It is irreducible,
since a reducible homogeneous quadratic is a product of two linear
forms and has rank at most two, whereas this one has rank six. On
the open set $p_{12}\ne0$, every tuple satisfying it is realized
by the columns
\[
 \binom10,\quad \binom0a,\quad
 \binom{-p_{23}/a}{p_{13}},\quad
 \binom{-p_{24}/a}{p_{14}},\qquad a=p_{12}.
\]
Thus this open set is a dense subset of the irreducible quadratic
hypersurface and there are no further polynomial relations. The
determinant $d$ is independently arbitrary, proving (P-F20f).

The Hilbert series and quotient are consequently
\[
 \operatorname{Hilb}_C(t)
   =\frac{1-t^4}{(1-t^2)^6(1-t^4)}
   =\frac1{(1-t^2)^6},\qquad
 C/(p_{ij}:i<j)=\mathbb C[d],\quad\deg d=4.
 \tag{P-F20g}
\]
The same exact-sequence argument as above forces $\Phi_4$ into
the universal denominator, although the actual Hilbert series is
regular at both primitive fourth roots.

We verify the stronger geometric conditions explicitly. The ring
$C$ is factorial. If a nonzero invariant polynomial is factored
into irreducibles in $\mathbb C[V]$, the connected group $G$
cannot permute its finite set of irreducible-factor lines
nontrivially. Each factor is therefore a semi-invariant, with
an algebraic character of $G$. Every character of $SL_2\times SL_4$
is trivial: the factors are generated by elementary unipotent
one-parameter subgroups, and an algebraic homomorphism from
$\mathbb G_a$ to $\mathbb G_m$ is constant because the only units
in $\mathbb C[t]$ are constants. Every irreducible factor is thus
itself invariant. Existence and uniqueness of factorization in
$C$ follow from those in $\mathbb C[V]$, and $C$ is normal.
As a hypersurface, it is Gorenstein and Cohen--Macaulay; see, for
example, the local-complete-intersection statement in the Stacks
Project, Tag~\texttt{0DWA}.

Rational singularities can also be seen directly in this example.
The quadratic factor is the affine cone over a smooth quadric
$Q^4\subset\mathbb P^5$. Blowing up its vertex gives the smooth
total space of $\mathcal O_{Q^4}(-1)$, with the zero section
contracted to the vertex. For every $k\ge0$, the sequence
\[
 0\longrightarrow\mathcal O_{\mathbb P^5}(k-2)
 \longrightarrow\mathcal O_{\mathbb P^5}(k)
 \longrightarrow\mathcal O_{Q^4}(k)\longrightarrow0
\]
and the projective-space cohomology formula (Stacks Project,
Tag~\texttt{01XT}) give
$H^i(Q^4,\mathcal O_{Q^4}(k))=0$ for every $i>0$.
The projection of this total space to $Q^4$ is affine, with
pushforward $\bigoplus_{k\ge0}\mathcal O_{Q^4}(k)$, so all its
higher structure-sheaf cohomology vanishes. The target cone is
affine; therefore the higher direct images under the resolution
vanish, and its direct image in degree zero is the structure
sheaf by normality. Taking the product of this resolution with
the $d$-line gives the same conclusion for $\operatorname{Spec}C$.

This counterexample rules out the proposed exact transfer even
within factorial Gorenstein rational invariant rings of connected
semisimple groups. It is a representation with several summands
and is not the irreducible cubic representation in P-F19. It does
not rule out a quantitative comparison of maximal denominator
weights; here the positive-form weight is two and the hidden
universal order is only four. The primary first-fundamental-theorem
source is
\url{https://dmi.unibas.ch/fileadmin/user_upload/dmi/Personen/Kraft_Hanspeter/Classical_Invariant_Theory.pdf}.

\paragraph{A precise obstruction to the proposed affine-slice argument.}
Large periods in affine slices do not, from multivariate
quasi-polynomiality alone, force large periods on rays through the
origin. Fix $M=2^a$ with $a\ge2$, and on
$\mathbb Z_{\ge0}^2$ define
\[
 F_M(x,y)=1+\mathbf1_{\{y\ {\rm odd},\ M\mid x\}}.
 \tag{P-F20e}
\]
This is a nonnegative multivariate quasi-polynomial, with
$F_M(0,0)=1$, and it has the rational generating series
\[
 \sum_{x,y\ge0}F_M(x,y)X^xY^y
  =\frac1{(1-X)(1-Y)}+
       \frac{Y}{(1-X^M)(1-Y^2)}.
\]
For every integral direction $(u,v)\in\mathbb Z_{\ge0}^2$,
the function $k\mapsto F_M(ku,kv)$ is either identically one
or is $1+\mathbf1_{\{k\ {\rm odd}\}}$, so its global period
divides two. Indeed, even $k$ makes $kv$ even, and odd $k$ is
invertible modulo $M$; hence for odd $k$ the extra term occurs
exactly when $v$ is odd and $M\mid u$.
But the fixed-target slice is
\[
 F_M(x,1)=1+\mathbf1_{\{M\mid x\}},
\]
whose generating series has a pole at a primitive $M$th root.
Thus cancellation among multivariate Fourier characters can hide
arbitrarily large powers of two on \emph{every} simultaneous ray.

This last example is not asserted to be the Hilbert function of
a branching domain. It pinpoints why a Fourier argument using
only a finite chamber decomposition and global
quasi-polynomiality on each ray cannot finish the transfer.
A successful argument must use additional structure of the actual
branching algebra. P-F20 leaves the weaker positive-form existence
clause and PH2 undecided; its new conclusions are the universal
factor and the actual, but fixed-target, covariant pole.
